\subsection{Sprint 1}

Após obter a apovação do protótipo, deu-se seguimento para o desenvolvimento real da aplicação. Para a sprint 1 estavam previstas as seguintes atividades: dedicação dos esforços da equipe em desenvolver o \textit{backend} do projeto na linguagem Java, previamente definida e o \textit{frontend} na linguagem Flutter; utilização do GIT, ambiente de controle de versões, bem como o uso da padronização estipulada, visando manter um código limpo e organizado; e comprometimento da equipe em desenvolver, finalizar e entregar as seguintes \textit{tasks}: criar "Tela de Cadastro de Atração Turística - Fotos", criar "Tela de Cadastro de Atração Turística - Horários"; criar "Tela de Cadastro de Atração Turística - Tags"; criar "Tela de Cadastro de Atração Turística - Categorias"; criar "Tela de Cadastro de Atração Turística - Formulário"; criar "Tela de Detalhes de Atração Turística"; criar "Tela de Listar Atrações Turísticas"; criar "Tela de Login"; criar "Rota de Listar uma Atração Turística por ID"; criar a rota de listagem de atrações turísticas; criar a rota de edição de atração turística; criar a rota de cadastro de atração turística; e criar a entidade atração turística e suas dependências.

Assim, durante o período da sprint, a equipe direcionou esforços para o desenvolvimento das tarefas estipuladas, visando cobrir algumas user stories contidas no escopo, de modo a entregar algumas telas ao cliente até o final do prazo estipulado. As \ac {us} atendidas durante essa sprint foram: \ac{us}01 - Cadastro de Atração Turística; \ac{us}02 - Listagem de Atrações Turísticas; \ac{us}03 - Visualização de Atração Turística; e \ac{us}04 - Tela de Login.

Até a data final da sprint, praticamente todas as tarefas foram concluídas com sucesso pela equipe, e as outras atividades previstas, anteriormente citadas, também foram finalizadas. Todos os integrantes da equipe cumpriram com suas obrigações e realizaram a entrega de suas respectivas \textit{tasks}, inclusive eu, que fiquei responsável pela produção da tela de Cadastro de Atração Turística - Categorias. 

Além disso, a dedicação da equipe em seguir as regras de \textit{commit} do projeto garantiu a organização da documentação e a iniciativa dos colegas em ajudar aqueles que precisavam fez uma diferença positiva no resultado final.

Entretanto, houveram alguns contratempos marcantes durante o andamento dessa print. A relativa falta de um melhor planejamento na delegação de tarefas e o exagero na quantidade de tarefas escolhidas para entrega dentro do prazo estipulado na sprint acabou sobrecarregando o time. Com isso, nem todos os integrantes conseguiram entregar suas respectivas tarefas dentro do prazo estipulado, deixando a US03 - Visualização de Atração Turística em débito para a próxima sprint.